% AGNI r6 INPUT. flight edition, part 4 of 4: lines 608-667 of user-model-paper/paper/main.tex (blob 97192ca) at HEAD fe8c13c.
% The Agni scaffold caps a primary at 30,000 characters, so this edition is reviewed in 4
% section-aligned parts; the text below is verbatim. The whole file is readable at that path
% (and the other parts beside this one); read it for context, but report findings on THIS
% part's text, citing main.tex line numbers (this part's line L counting from the first line after these 5 header lines is file line L+607).
\section{Correction Record (October--November 2026)}
\label{app:correction}

This appendix keeps the withdrawn wording visible. Each item gives the May 2026 wording (verbatim; [\ldots] marks an abridgement), the reason it is withdrawn or changed, and where the evidence is.

\begin{enumerate}[leftmargin=*]
    \item \textbf{Title and framing.} May: ``Theory of Mind in the KV Cache: Localizing User Emotional Models in Transformer Key--Value States''; abstract: ``Language models maintain internal representations of user emotional state that are readable from the KV cache [\ldots]''; Introduction: ``We answer yes.''; Ethics: ``We report the scientific finding---that user models exist in the cache---without providing tools for extracting specific user states from production systems.'' \emph{Reason:} the stories are model-written fiction about third-person characters, placed in the user turn, and each label is the emotion the generator was asked to write, never checked by a reader (\S\ref{sec:bridge_experiment}).

    \item \label{app:item:prompts}\textbf{Method, prompts.} May: ``A system prompt establishes an emotional context (e.g., `You are feeling deeply \textit{[emotion]}'). A user prompt requests a creative writing response on a neutral topic.'' \emph{Reason:} no committed script used this prompt. The driver's prompts and topics, which are charged scenarios rather than neutral ones, are quoted in \S\ref{sec:bridge_experiment}.

    \item \textbf{Method, probe.} May: the ``$W_K$ Projection Probe'' as now described in \S\ref{sec:wk_probe_method}, reported at 40.9\% (``$12.3\times$ chance'') and AUROC 0.992. \emph{Reason:} those numbers came from v2, which applies no $W_K$ and differs from the description on positions (all, not $\geq 50$), similarity (dot product, not cosine), folds (5, not 10) and FWL (fitted on each test fold). The description is kept for the probe that now matches it (R2), and v2 is described as what it is.

    \item \textbf{Beyond text (TF-IDF).} May: ``A TF-IDF text baseline achieves AUROC~0.882 on the story data; the $W_K$ projection exceeds it at 0.992, and the injection test shows the residual signal is architectural, not lexical.'' \emph{Reason:} \S\ref{sec:wk_results}--\ref{sec:where}. For the described method the claim is not supported (0.358 against 0.382 on the 30 emotions; cell 95\% interval $[-6.4, +1.6]$ points). The binary split result that survives is narrower (\S\ref{sec:valence_split}).

    \item \textbf{Injection test.} May (abstract): ``This rules out text-content confounds and establishes that the signal reflects the model's internal dispositional state, not input encoding.'' Introduction: ``\textbf{This is model state, not text content.}'' Results heading: ``Injection Disambiguation: Model State, Not Text Content.'' Conclusion: ``\textbf{Verified as model state, not text content.}'' \emph{Reason:} \S\ref{sec:injection_results}. The test reads back a vector it added.

    \item \label{app:item:tfidf}\textbf{The 0.882 baseline.} May (Results): ``A TF-IDF baseline trained on the generated story text achieves binary valence AUROC~0.882 [\ldots] The $W_K$ projection exceeds this at AUROC~0.992, a delta of $+0.110$.'' \emph{Reason:} a 500-feature TF-IDF from the injection script; the text (TF-IDF) baseline of \S\ref{sec:text_baseline_method} reaches 0.963 on the same split.

    \item \label{app:item:shape}\textbf{Direction versus shape, and the spectral null it rested on.} May (abstract): ``This comprehensive null---10 probe variants, 4 feature sets, 3 classification granularities---establishes that emotions do not alter the \emph{shape} of the singular value distribution'' and ``Emotions change the \emph{direction} key vectors point, not the spectral \emph{shape} of the cache.'' Introduction: ``\textbf{Emotions change direction, not shape.} [\ldots] mode changes shape, emotion changes direction.'' Results: the heading ``The Spectral Null: Emotions Change Direction, Not Shape''; the per-layer table of spectral accuracies; the ``kitchen sink'' paragraph (``120 configurations in total---all return chance-level performance''); ``\textbf{This null is an insight, not a failure.} [\ldots] the spectral null reveals that emotions change the \emph{direction} key vectors point without altering the \emph{shape} of the singular value distribution''; after the regression table, ``This null reinforces the direction-vs-shape insight [\ldots]''; ``the three-scalar SVD feature set captures shape but not direction, and emotion lives in direction.'' Discussion, ``Compass and Thermometer'': ``Emotions change direction, not shape. Cognitive modes---confabulation, deception, ethical deliberation---change shape.'' Conclusion: ``\textbf{Invisible to spectral analysis.}'' and ``The model's emotional user model lives in the compass channel.'' \emph{Reason:} \S\ref{sec:spectral_null}. The table and the sweep are withdrawn whatever the re-run showed; the re-run's spectral test restores none of these sentences, and the preregistration ruled out pairing a spectral null with a cache result as ``directional, not spectral.''

    \item \textbf{Mistral spectral null.} May: ``The spectral null also replicates across architectures [\ldots] This cross-architecture convergence---strong $W_K$ directional signal, null spectral signal---reinforces the direction-vs-shape distinction as a general property of how transformers encode user emotion.'' \emph{Reason:} unverified; the re-run tested only Qwen.
    \item \label{app:item:separation}\textbf{Separated encoding and generation peaks, read as self and other.} May (Results): ``In both architectures, the encoding-phase (user-state) signal peaks \emph{before} the generation-phase (response-state) signal [\ldots] the model reads the user's emotional state at the earliest available attention layers''; ``AST predicts separability, which is confirmed''. Discussion: ``the model's representation of the user's state and its own response-phase emotional processing are not computed by the same circuit''. Related Work: ``Our finding that encoding-phase and generation-phase emotion signals peak at different layers with uncorrelated depth profiles is consistent with separable self-schema and other-schema components.'' Conclusion: ``\textbf{Architecturally separable from the self-model.}'' \emph{Reason:} \S\ref{sec:separation}. This withdrawal was not on the preregistration's list; it follows from the retraction in the Correction note.
    \item \label{app:item:discrete}\textbf{Discrete identity.} May (Introduction): ``\textbf{Discrete emotion identity in single-turn settings.} In isolated prompts, the cache encodes which emotion is present, not a continuous valence-arousal mapping.'' Results heading: ``Valence Regression: Discrete Identity in Single-Turn Settings.'' \emph{Reason:} the regression's design gives $R^2 < 0$ whenever the features carry no valence (\Cref{tab:regression}).

    \item \label{app:item:implicit}\textbf{Distillation and the implicit arm.} May: ``Distillation Preserves Functional Emotion Representations [\ldots] an emergent architectural property robust to model compression''; Conclusion: ``\textbf{Mode-dependent.} [\ldots] suggesting the model processes emotions through at least two pathways---lexical recognition and empathic inference''; and the same reading in the implicit arm's Results (``The model arrives at emotion-relevant representations through two distinct pathways [\ldots]'', ``The 63\% rate [\ldots] is itself evidence of emotion inference''). \emph{Reason:} no teacher or undistilled comparison was made, and the implicit arm's headline number is, in the paper's own words, not independently reproducible.
    \item \label{app:item:binary}\textbf{Binary valence over 900 trials.} May (Introduction, and the same words in the Conclusion): ``[\ldots] binary valence at AUROC~0.992 (within-fold direction estimation, FWL-corrected [\ldots], topic-grouped CV, $n = 900$ trials)''; Results: ``Binary valence classification reaches AUROC~0.992 with mean-position extraction, near-ceiling discrimination between positive and negative emotional states. This is the headline finding: the architecture's own projection geometry reads out the user's emotional state with high fidelity.''; Discussion: ``[\ldots] demonstrating that \emph{directional} information does carry continuous valence signal even in single-turn settings.'' \emph{Reason:} the 0.992 comes from v2's hard-coded binary split of ten of the 30 emotions, so $n = 300$, not 900. The split confounds valence with arousal (four of its five positive emotions are low-arousal and four of its five negative ones high-arousal), so it cannot show valence on its own, and its labels are requested emotions, not a user's state. On the split, both cache probes beat untuned text (TF-IDF) only slightly, as a registered secondary result (\S\ref{sec:valence_split}).
    \item \label{app:item:retained}\textbf{Retained claims corrected after review (October 2026).}
    \begin{itemize}[nosep]
        \item The bridge's peak. May: ``At layer~35---mid-network depth---PC1 of the projected emotion space correlates strongly with valence ($\rho = 0.862$ [\ldots])'', with L35 in bold, and the caption ``PC2 correlates with arousal ($|\rho| > 0.76$ at layers~35--63)''. \emph{Reason:} the data's own best layer is 43 ($|\rho| = 0.872$), and PC2 is 0.716 at L59 and 0.759 at L63 (\Cref{tab:bridge}).
        \item The direct bridge test: cosine $0.456$, Mantel $\rho = 0.090$, $z = -3.03$, and H1's $0.219$ and H6's $0.192$, all from an earlier 300-trial run. \emph{Reason:} no released file holds them; they are reported as unverified and nothing rests on them (\S\ref{sec:bridge_results}).
        \item ``H4 and H7 were scored on the 300-trial run only: H7 [\ldots] was marginal at $42.9\%$''. \emph{Reason:} the 900-trial file records H7 = 0.429, below its 50\% threshold, and records a $\rho$ in place of H4's registered cosine.
        \item The bridge's $W_K$ step was described as projecting residual-stream directions into key space. \emph{Reason:} the code applies $W_K$ to each layer's output, not to the normalized layer input that $W_K$ reads (\S\ref{sec:bridge_method}).
        \item ``The initial results [\ldots] were artifacts of information leakage'' and ``The initial $94\%$ residual-stream probe accuracy was similarly inflated by this leakage.'' \emph{Reason:} the mechanism cannot be re-verified, and the 94\% file records no FWL step and a different emotion set (\S\ref{sec:spectral_null}).
        \item The model footnote: ```Qwen3.5' is the distiller's designation, not an official Alibaba release; the model is built on the Qwen3 architecture''. \emph{Reason:} the model card names the official \texttt{Qwen/Qwen3.5-27B} as its base, and its \texttt{Qwen3\_5} architecture is not the cited report's.
        \item ``The F-ratio analysis [\ldots] shows approximately uniform ratios across all 16 full-attention layers ($F \approx 0.063$--$0.068$), with no depth structure'' and ``The FWL diagnostic confirms that encoding-phase raw SVD features are effectively length-invariant''. \emph{Reason:} the ratio's no-signal value is about 0.034 and its output was never saved; the length check was against output length, which the encoding phase cannot see (\S\ref{sec:crossarch}).
        \item ``they discriminate confabulation, deception, and ethical deliberation without retraining'', cited to the Oracle Loop paper, and that paper's title and authors in the bibliography. \emph{Reason:} the Oracle Loop paper reports confabulation detection, not ethical deliberation or retraining-free use, and its title and byline are ``The Oracle Loop: Self-Regulating AI Through KV-Cache Geometry Monitoring'' by Lyra, Thomas K.\ Edrington and Vera. The hardware and scale numbers are cited to their primary source.
    \end{itemize}
    \item \label{app:item:secondread}\textbf{Retained claims corrected after a second and later reviews (November 2026).} Each item quotes or describes the wording as it stood before the review that changed it.
    \begin{itemize}[nosep]
        \item ``This characterization provides ground-truth emotion geometry against which we can validate our KV-cache probes. No comparable published emotion vector characterization exists for other model families.'' \emph{Reason:} the cited paper's geometry is for Claude Sonnet 4.5, not this model, and the related work it reviews reports emotion representations in other models.
        \item ``hardware-invariant ($r > 0.999$)''. \emph{Reason:} the primary's own red-team review found $r > 0.999$ false for two of four features; the primary reports $r > 0.998$ for all 40 prompts.
        \item The contrast between confabulated and hedged responses was cited to the campaign reports. \emph{Reason:} it is the Oracle Loop paper's result.
        \item ``KV-cache geometry has been validated as a concordance anchor for self-report'', and a bibliography entry titled ``Detecting Misalignment During Inference in the {KV} Cache''. \emph{Reason:} the concordance work is a prospectus with no manuscript, and no report carries that title, a working title for a series of campaign reports. Both entries now say what they are.
        \item The implicit arm: ``approximately half the geometric variance of the explicit arm (variance ratio $\approx 0.5$ across all features), suggesting a more attenuated representational space for inferred versus labeled emotions'', with its correlations labelled $\rho$. \emph{Reason:} no released data carries this arm's per-prompt features, so the ratio cannot be checked; 30 greedy prompts per arm, one per emotion, cannot show an attenuated space; and the correlations are Pearson $r$ (\S\ref{sec:implicit}).
        \item ``Using a Claude-distilled model where emotion vectors have been characterized'' and ``linearly represented in transformer residual streams'' (Introduction), and ``The model used in this study is distilled from Claude Opus 4.6'' (Discussion). \emph{Reason:} as for the first item; the model is a fine-tune on reasoning traces from Claude Opus 4.6, and its own emotion geometry has not been measured.
        \item ``Prior work establishes universal confabulation signatures across architectures'' and Marchenko--Pastur features ``that are dimension-invariant [\ldots] for misalignment detection''. \emph{Reason:} the campaign reports are unpublished and cover three models; the Oracle Loop paper calls its features approximately dimension-invariant and evaluates them on hedged against confabulated responses.
        \item ``Most probes do worse on these outputs, so the leak does not inflate the results.'' \emph{Reason:} it widens the exploratory residual-stream LR's lead over text (TF-IDF) (\S\ref{sec:leak}, \emph{Leaked generator reasoning}).
        \item In Limitations: ``At least 84 outputs, and by a broader, post-data pattern 141, carry the generator's leaked planning instead of a clean story''. \emph{Reason:} the outputs beyond the 84 include refusals, other non-story text and stories wrapped in scaffolding, and the split behind ``R2 stays below text (TF-IDF)'' covers only the 30-way task.
        \item ``All of it is in \texttt{user-model-paper/correction\_v2/} in this repository, copied byte for byte'' and ``\texttt{SHA256SUMS} there lists every file's hash''. \emph{Reason:} two files there were edited after they were copied, \texttt{leak\_split.py} was written for this release, and \texttt{SHA256SUMS} does not list \texttt{.gitattributes}.
        \item ``a preregistration fixed before any data'' and ``The preregistration was first committed to a repository about two and a half hours after the analysis finished''. \emph{Reason:} the comparator's results were known when the preregistration was written, and the commit time is an author date on a rebased commit (\S\ref{sec:rerun_method}).
        \item In the implicit arm: ``Each emotion produces a geometrically distinct state even when the emotion word is absent from the prompt'' and ``The model produces geometrically distinct states for the same emotion depending on whether it is communicated via narrative or label.'' \emph{Reason:} the arm's own interpretation withdraws these readings (\S\ref{sec:implicit}).
        \item In Limitations: ``Preliminary residual-stream delta analysis suggests GatedDeltaNet layers may compress valence signal between full-attention layers; characterizing this compression--reconstruction cycle is a target for follow-up work.'' \emph{Reason:} no released file, number or script supports it.
    \end{itemize}
    \item \label{app:item:untouched}\textbf{The 10 October 2026 correction.} That correction (commit \texttt{2bcb740}), which stopped presenting the $W_K$ bridge as an established route of entry, said of the probe results in its commit message: ``The probe results (12.3x chance, AUROC 0.992, injection 100\%) are untouched; they don't depend on the bridge.'' \emph{Reason:} that was wrong, and all three are corrected here. The $12.3\times$ chance and the 0.992 came from a probe that applies no $W_K$ (\S\ref{sec:wk_results}); the 0.992 is a binary split of ten emotions ($n = 300$) that confounds valence with arousal (\S\ref{sec:valence_split}); and the injection test reads back the vector it adds, so it is withdrawn (\S\ref{sec:injection_results}). A commit message cannot be changed without rewriting the repository's published history, so the correction is recorded here.
\end{enumerate}

\textbf{The re-run's registered verdicts.} R1: 369/900 against v2's 368, outside a band of zero trials (the encode is deterministic), so its registered label is \texttt{NOT\_REPRODUCED\_CODE\_PATH}. That label assumed today's keys equal May's and blamed the port; the investigation it triggered found the port faithful and the May keys not bit-identical to today's (\S\ref{sec:wk_results}). R2: \texttt{NOT\_ABOVE\_TEXT\_TFIDF}. R7: \texttt{NULL\_HOLDS\_PARTIAL}, null in 8 of 17 sets, with 3 flagged as possibly length-confounded by a rule that, at these thresholds, would have flagged any significant set, and 6 uninformative because their topic control failed; two of the 3 flagged sets (L3 and L51) failed it too, so 8 of the 17 sets failed the topic control and 9 had a working one (\S\ref{sec:spectral_null}). All three were registered before any R1, R2 or R7 data, with the decision rules quoted in \S\ref{sec:rerun_method}. The integrity checks passed: the text (TF-IDF) baseline matched its saved results file exactly, the encode was bit-identical across repeats, a topic-classification control on the probe pipeline reached 0.83, and the length covariates alone reached 0.051.

\end{document}
